\newpage
\phantomsection
\label{prf:fv-msml-distance-to-set-eq-zero-of-mem}
\LRAProofFor{thm:fv-msml-distance-to-set-eq-zero-of-mem}

\begin{remark*}[Return]
Return to \hyperref[thm:fv-msml-distance-to-set-eq-zero-of-mem]{Distance to a Set Vanishes at a Point of the Set}.
\end{remark*}

\begin{theorem*}[Distance to a Set Vanishes at a Point of the Set]
Let $(X,d)$ be a metric space. If \(x\in A\), then
\[
  d(x,A)=0.
\]
\end{theorem*}

\begin{proof}[Professional Standard Proof]
\LRAProofBodyStart
First, because \(x\in A\), the witness estimate gives
\[
  d(x,A)\leq d(x,x)=0.
\]
Second, every distance in the set \(\{d(x,y)\mid y\in A\}\) is nonnegative, so
\(0\) is a lower bound for that distance set. Since \(d(x,A)\) is the greatest
lower bound, \(0\leq d(x,A)\). Combining \(d(x,A)\leq0\) and
\(0\leq d(x,A)\) gives \(d(x,A)=0\).
\end{proof}

\begin{proof}[Detailed Learning Proof]
\LRAProofBodyStart
The Lean proof proves the equality by proving both inequalities and then using
\texttt{le\_antisymm}.

For the lower bound \(0\leq d(x,A)\), Lean first constructs nonemptiness of
\(A\) from the hypothesis \(x\in A\). The witness is \(x\) itself:
\[
  \langle x,\ \texttt{point\_in\_set}\rangle.
\]
With this witness, \texttt{distanceToSet\_isGLB} gives
\[
  d(x,A)\ \text{is the greatest lower bound of}\ D_{x,A},
\]
where \(D_{x,A}\) denotes the Lean set \texttt{distanceSet x A}.

Next Lean proves that \(0\) is a lower bound for the distance set. To do this,
take an arbitrary real number \(r\in D_{x,A}\). By the definition of
\texttt{distanceSet}, this membership has a witness
\(y\in A\) such that
\[
  r=d(x,y).
\]
After rewriting \(r\) by this equality, the required inequality becomes
\[
  0\leq d(x,y),
\]
which is exactly nonnegativity of metric distance. This is the Lean step
\texttt{rcases distance\_value\_in\_set with <y, \_point\_in\_set, rfl>}
followed by \texttt{dist\_nonneg}.

Since \(0\) is a lower bound and \(d(x,A)\) is the greatest lower bound, the
greatest-lower-bound property gives
\[
  0\leq d(x,A).
\]

For the upper bound \(d(x,A)\leq 0\), Lean uses \(x\) itself as the witness
point in \(A\). The previous theorem gives
\[
  d(x,A)\leq d(x,x)=0.
\]
The first inequality is
\texttt{distanceToSet\_le\_distance\_to\_point\_of\_mem x point\_in\_set};
the equality \(d(x,x)=0\) is \texttt{dist\_self x}. These are chained in a
\texttt{calc} block.

Finally, Lean combines
\[
  d(x,A)\leq 0
  \qquad\text{and}\qquad
  0\leq d(x,A)
\]
with \texttt{le\_antisymm}, producing \(d(x,A)=0\).
\end{proof}

\begin{proof}[Formalized Lean Proof]
\begin{verbatim}
theorem distanceToSet_eq_zero_of_mem
    {X : Type u}
    [MetricSpace X]
    {A : Set X}
    {x : X}
    (point_in_set : x in A) :
    distanceToSet x A = 0 := by
  have distanceToSet_nonnegative : 0 <= distanceToSet x A := by
    have A_nonempty : A.Nonempty := <x, point_in_set>
    have infimum_property :
        IsGLB (distanceSet x A) (distanceToSet x A) := by
      exact distanceToSet_isGLB x A_nonempty
    have zero_lower_bound : 0 in lowerBounds (distanceSet x A) := by
      intro r distance_value_in_set
      rcases distance_value_in_set with <y, _point_in_set, rfl>
      exact dist_nonneg
    exact infimum_property.2 zero_lower_bound
  have distanceToSet_nonpositive : distanceToSet x A <= 0 := by
    calc
      distanceToSet x A <= dist x x := by
        exact distanceToSet_le_distance_to_point_of_mem x point_in_set
      _ = 0 := by
        exact dist_self x
  exact le_antisymm distanceToSet_nonpositive distanceToSet_nonnegative
\end{verbatim}
\end{proof}

\begin{remark*}[Proof structure]
The final line uses \texttt{le\_antisymm}: equality in an ordered type is
proved by proving both inequalities.
\end{remark*}

\NoLocalDependencies

\clearpage
